%% Beginning of file 'sample701.tex'
%%
%% Version 7.0.1. Created May 2025.
%% Version 7. Created January 2025.  
%%
%% AASTeX v7+ calls the following external packages:
%% times, hyperref, ifthen, hyphens, longtable, xcolor, 
%% bookmarks, array, rotating, ulem, and lineno 
%%
%% RevTeX is no longer used in AASTeX v7+.
%%
\documentclass[trackchanges, twocolappendix, twocolumn]{aastex701}
%%
%% This initial command takes arguments that can be used to easily modify 
%% the output of the compiled manuscript. Any combination of arguments can be 
%% invoked like this:
%%
%% \documentclass[argument1,argument2,argument3,...]{aastex701}
%%
%% Six of the arguments are typestting options. They are:
%%
%%  twocolumn   : two text columns, 10 point font, single spaced article.
%%                This is the most compact and represent the final published
%%                derived PDF copy of the accepted manuscript from the publisher
%%  default     : one text column, 10 point font, single spaced (default).
%%  manuscript  : one text column, 12 point font, double spaced article.
%%  preprint    : one text column, 12 point font, single spaced article.  
%%  preprint2   : two text columns, 12 point font, single spaced article.
%%  modern      : a stylish, single text column, 12 point font, article with
%% 		  wider left and right margins. This uses the Daniel
%% 		  Foreman-Mackey and David Hogg design.
%%
%% Note that you can submit to the AAS Journals in any of these 6 styles.
%%
%% There are other optional arguments one can invoke to allow other stylistic
%% actions. The available options are:
%%
%%   astrosymb    : Loads Astrosymb font and define \astrocommands. 
%%   tighten      : Makes baselineskip slightly smaller, only works with 
%%                  the twocolumn substyle.
%%   times        : uses times font instead of the default.
%%   linenumbers  : turn on linenumbering. Note this is mandatory for AAS
%%                  Journal submissions and revisions.
%%   trackchanges : Shows added text in bold.
%%   longauthor   : Do not use the more compressed footnote style (default) for 
%%                  the author/collaboration/affiliations. Instead print all
%%                  affiliation information after each name. Creates a much 
%%                  longer author list but may be desirable for short 
%%                  author papers.
%% twocolappendix : make 2 column appendix.
%%   anonymous    : Do not show the authors, affiliations, acknowledgments,
%%                  and author contributions for dual anonymous review.
%%  resetfootnote : Reset footnotes to 1 in the body of the manuscript.
%%                  Useful when there are a lot of authors and affiliations
%%		    in the front matter.
%%   longbib      : Print article titles in the references. This option
%% 		    is mandatory for PSJ manuscripts.
%%
%% Since v6, AASTeX has included \hyperref support. While we have built in 
%% specific %% defaults into the classfile you can manually override them 
%% with the \hypersetup command. For example,
%%
%% \hypersetup{linkcolor=red,citecolor=green,filecolor=cyan,urlcolor=magenta}
%%
%% will change the color of the internal links to red, the links to the
%% bibliography to green, the file links to cyan, and the external links to
%% magenta. Additional information on \hyperref options can be found here:
%% https://www.tug.org/applications/hyperref/manual.html#x1-40003
%%
%% The "bookmarks" has been changed to "true" in hyperref
%% to improve the accessibility of the compiled pdf file.
%%
%% If you want to create your own macros, you can do so
%% using \newcommand. Your macros should appear before
%% the \begin{document} command.
%%
\usepackage{amsmath}
\usepackage{bm}
\usepackage{booktabs}
\usepackage{CJK}
\usepackage{comment}
\usepackage{graphicx}
\usepackage{makecell}
\usepackage{multirow} % multiple rows
\usepackage{subcaption}
\usepackage{txfonts} % Times New Roman

\newcommand{\vdag}{(v)^\dagger}
\newcommand\aastex{AAS\TeX}
\newcommand\latex{La\TeX}

\newcommand{\czkong}[1]{{\color{magenta}[\textbf{#1}]}}
%%%%%%%%%%%%%%%%%%%%%%%%%%%%%%%%%%%%%%%%%%%%%%%%%%%%%%%%%%%%%%%%%%%%%%%%%%%%%%%%
%%
%% The following section outlines numerous optional output that
%% can be displayed in the front matter or as running meta-data.
%%
%% Running header information. A short title on odd pages and 
%% short author list on even pages. Note that this
%% information may be modified in production.
\shorttitle{\textsc{Arepo-CD}}
\shortauthors{C. Kong et al.}
%%
%% Include dates for submitted, revised, and accepted.
\received{\today}
\revised{\today}
\accepted{\today}
%%
%% Indicate AAS Journal the manuscript was submitted to.
\submitjournal{APJL}
%% Note that this command adds "Submitted to " the argument.
%%
%% You can add a light gray and diagonal water-mark to the first page 
%% with this command:
%% \watermark{text}
%% where "text", e.g. DRAFT, is the text to appear.  If the text is 
%% long you can control the water-mark size with:
%% \setwatermarkfontsize{dimension}
%% where dimension is any recognized LaTeX dimension, e.g. pt, in, etc.
%%%%%%%%%%%%%%%%%%%%%%%%%%%%%%%%%%%%%%%%%%%%%%%%%%%%%%%%%%%%%%%%%%%%%%%%%%%%%%%%
%%
%% Use this command to indicate a subdirectory where figures are located.
%%\graphicspath{{./}{figures/}}
%% This is the end of the preamble.  Indicate the beginning of the
%% manuscript itself with \begin{document}.

\begin{document}
\begin{CJK*}{UTF8}{gbsn}

\title{Connecting Intermediate-Mass Black Holes Formed in Proto-Globular Clusters to High Redshift Supermassive Black Hole Seeds}

%% A significant change from AASTeX v6+ is in the author blocks. Now an email
%% address is required for each author. This means that each author requires
%% at least one of the following:
%%
%% \author
%% \affiliation
%% \email
%%
%% If these three commands are not available for each author, the latex
%% compiler will issue an error and if you force the latex compiler to continue,
%% it will generate an incomplete pdf.
%%
%% Multiple \affiliation commands are allowed and authors can also include
%% an optional \altaffiliation to indicate a status, i.e. Hubble Fellow. 
%% while affiliations are indexed as footnotes, altaffiliations are noted with
%% with a non-numeric footnote that is set away from the numeric \affiliation 
%% footnotes. NOTE that if an \altaffiliation command is used it must 
%% come BEFORE the \affiliation call, right after the \author command, in 
%% order to place the footnotes in the proper location. Because non-numeric
%% symbols are used, \altaffiliation should be used sparingly.
%%
%% In v7+ the \author command takes an optional argument which provides 
%% additional metadata about the author. Authors can provide the 16 digit 
%% ORCID, the surname (family or last) name, the given (first or fore-) name, 
%% and a name suffix, e.g. "Jr.". The syntax is:
%%
%% \author[orcid=0000-0002-9072-1121,gname=Gregory,sname=Schwarz]{Greg Schwarz}
%%
%% This name metadata in not shown, it is only for parsing by the peer review
%% system so authors can be more easily identified. This name information will
%% also be sent to the publisher so they can include it in the CROSSREF 
%% metadata. Including an orcid will hyperlink the author name to the 
%% author's ORCID page. Note that  during compilation, LaTeX will do some 
%% limited checking of the format of the ID to make sure it is valid. If 
%% the "orcid-ID.png" image file is  present or in the LaTeX pathway, the 
%% ORCID icon will appear next to the authors name.
%%
%% Even though emails are now required for each author, the \email does not
%% produce output in the compiled manuscript unless the optional "show" command
%% is used. For example,
%%
%% \email[show]{greg.schwarz@aas.org}
%%
%% All "shown" emails are show in the bottom left of the first page. Due to
%% space constraints, only a few emails should be shown. 
%%
%% To identify a corresponding author, use the \correspondingauthor command.
%% The command appends "Corresponding Author: " to the argument it appears at
%% the bottom left of the first page like the output from \email. 

\author[0009-0003-9099-5805,gname=Chuizheng,sname=Kong]{Chuizheng Kong (孔垂政)}
\affiliation{Department of Astronomy, Tsinghua Univeristy, Beijing, 100084, People's Republic of China}
\email{kcz24@mails.tsinghua.edu.cn}

\author[0000-0002-1253-2763,gname=Hui,sname=Li]{Hui Li (李辉)}
\affiliation{Department of Astronomy, Tsinghua Univeristy, Beijing, 100084, People's Republic of China}
\email[show]{hliastro@tsinghua.edu.cn}

%% Use the \collaboration command to identify collaborations. This command
%% takes an optional argument that is either a number or the word "all"
%% which tells the compiler how many of the authors above the command to
%% show. For example "\collaboration[all]{(DELVE Collaboration)}" wil include
%% all the authors above this command.
%%
%% Mark off the abstract in the ``abstract'' environment. 
\begin{abstract}

We describe the current semi-analytic nuclear-star-cluster (NSC) model used to connect globular cluster (GC) formation, GC orbital evolution, and intermediate-mass black hole (IMBH) seeding. The model paints GC formation onto cosmological halo merger trees, samples individual clusters from a Schechter initial cluster mass function, assigns birth radii from a Gao+2023 S{\'e}rsic spatial model, and evolves the clusters in an analytic, time-dependent galactic background. Stellar mass is removed by stellar evolution, continuous tidal stripping, direct tidal tearing, and dynamical friction-driven sinking. IMBHs are assigned once at cluster birth using a metallicity- and surface-density-dependent runaway-collision fit. If the stellar component of an IMBH-hosting GC is removed before the IMBH reaches the centre, the black hole is evolved as a wandering remnant until it either sinks or remains in the galaxy.

\end{abstract}

%% Keywords should appear after the \end{abstract} command. 
%% The AAS Journals now uses Unified Astronomy Thesaurus (UAT) concepts:
%% https://astrothesaurus.org
%% You will be asked to selected these concepts during the submission process
%% but this old "keyword" functionality is maintained in case authors want
%% to include these concepts in their preprints.
%%
%% You can use the \uat command to link your UAT concepts back its source.
\keywords{\uat{Galaxies}{573} --- \uat{Cosmology}{343} --- \uat{High Energy astrophysics}{739} --- \uat{Interstellar medium}{847} --- \uat{Stellar astronomy}{1583} --- \uat{Solar physics}{1476}}

%% From the front matter, we move on to the body of the paper.
%% Sections are demarcated by \section and \subsection, respectively.
%% Observe the use of the LaTeX \label
%% command after the \subsection to give a symbolic KEY to the
%% subsection for cross-referencing in a \ref command.
%% You can use LaTeX's \ref and \label commands to keep track of
%% cross-references to sections, equations, tables, and figures.
%% That way, if you change the order of any elements, LaTeX will
%% automatically renumber them.

\section{Introduction} \label{sect:intro}

The discovery of massive black holes in the first billion years of cosmic history provides a direct constraint on the formation of black-hole seeds and their early assembly. Luminous quasars at $z \gtrsim 6$ host black holes with masses of $M_{\bullet} \sim 10^8$--$10^{10}~M_\odot$, requiring either massive initial seeds, sustained near-Eddington accretion, episodes of super-Eddington growth, low radiative efficiencies, or a combination of these processes \citep[e.g.,][]{Banados+2018SMBH, Matsuoka+2023SMBH, Larson+2023SMBH, Kocevski+2023SMBH, Harikane+2023SMBH, Bogdan+2024SMBH, Furtak+2024SMBH}. The tension has become sharper with JWST, which has revealed a population of fainter active galactic nuclei at high redshift, including systems whose inferred black-hole masses appear large compared with the stellar masses of their young host galaxies \citep[e.g.,][]{Maiolino+2024GNz11, Natarajan+2024UHZ1, Furtak+2024A2744QSO1}. These objects extend the early-black-hole problem from the most luminous quasars to a broader population of lower-mass galaxies, where the ratio $M_{\bullet}/M_\star$ may be even more informative than $M_{\bullet}$ alone.

A particularly important class of objects is the population of high-redshift black holes that appear overmassive relative to their hosts. In this work, we use ``overmassive'' in the observational sense, referring to systems whose inferred $M_{\bullet}/M_\star$ lies substantially above the local black-hole--galaxy scaling relations \citep[e.g.,][]{KormendyHo2013, ReinesVolonteri2015}. A2744-QSO1 provides a clear example: deep JWST/NIRSpec spectroscopy identifies it as a reddened broad-line AGN at $z \simeq 7.05$, with a black-hole mass of order $10^7~M_\odot$ and a black-hole-to-galaxy mass ratio well above that of typical local galaxies \citep{Furtak+2024A2744QSO1}. Similar high ratios have been inferred for other JWST-selected broad-line AGN, although the quantitative values remain sensitive to virial mass estimators, bolometric corrections, host-galaxy decomposition, obscuration, and selection effects \citep{Pacucci+2023, Harikane+2023, Greene+2024}. The relevant theoretical question is therefore not only how to form the most massive $z \gtrsim 6$ quasars, but also how to produce a high-redshift population with a broad and possibly biased distribution of $M_{\bullet}/M_\star$.

This question is closely connected to the nature of little red dots (LRDs). LRDs are compact red sources identified by JWST, often showing blue rest-UV continua, red rest-optical continua, strong Balmer breaks or Balmer-limit features, and in many cases broad Balmer emission lines \citep{Labbe+2023, Matthee+2024, Kocevski+2024, Greene+2024}. These properties suggest that at least a subset of LRDs host accreting black holes. At the same time, their spectral-energy distributions are difficult to interpret uniquely: stellar populations, dust-reddened AGN continua, dense gas around the accretion flow, and scattered or reprocessed AGN light can all contribute to the observed colours and line features \citep{Setton+2024, InayoshiMaiolino2025, Brooks+2024}. We therefore do not attempt to model the full LRD phenomenon. Instead, we focus on the subset of LRDs with evidence for black-hole components, asking whether their apparently large $M_{\bullet}/M_\star$ ratios can arise naturally from early seed formation and hierarchical assembly in compact stellar systems.

Most theoretical models of early black-hole formation are commonly grouped into light-seed and heavy-seed scenarios. Light seeds, produced by Population III stellar remnants, begin with masses of order $10$--$10^2~M_\odot$ and require efficient subsequent growth to explain the earliest massive black holes. Heavy seeds, produced by direct collapse or supermassive-star collapse in special atomic-cooling environments, may begin with $10^4$--$10^6~M_\odot$ and therefore reduce the required amount of later accretion \citep{BrommLoeb2003, Begelman+2006, LodatoNatarajan2006, LatifFerrara2016}. A third route lies between these limits: dense stellar systems can form intermediate-mass black holes through runaway stellar collisions, very massive stars, compact-remnant interactions, or repeated mergers in young massive clusters, globular-cluster progenitors, and nuclear star clusters \citep{PortegiesZwartMcMillan2002, DevecchiVolonteri2009, Davies+2011, Rantala+2026}. This channel is especially relevant at high redshift, where gas-rich galaxies are expected to form compact stellar systems efficiently.

Nuclear star clusters provide a natural mechanism for connecting cluster-scale seed formation to galaxy-scale black-hole assembly. Observationally, NSCs are dense stellar systems at galactic centres, and their formation is usually interpreted as a combination of in-situ star formation and the inspiral of massive star clusters through dynamical friction \citep{Neumayer+2020}. In a hierarchical universe, the same process can occur not only along the main progenitor branch, but also inside satellite galaxies before they merge into the main system. Satellite NSCs can therefore deliver both stellar mass and embedded black-hole seeds to the growing central nucleus. This channel is particularly relevant for overmassive high-redshift black holes because it can concentrate black-hole mass into small galaxies before their stellar hosts have fully assembled.

In this paper, we develop a semi-analytic model for this compact-stellar-system pathway. GC formation is painted onto cosmological merger trees and tied to halo growth, gas availability, a Schechter cluster initial mass function, and spatial sampling in host or satellite systems; IMBH seeds are then assigned at cluster birth using the metallicity- and surface-density-dependent runaway-collision fit from FROST-CLUSTERS III \citep{Rantala+2026}. The subsequent evolution follows stellar mass loss, tidal stripping, direct tidal tearing, and dynamical-friction sinking, allowing intact clusters, disrupted remnants, wandering IMBHs, and satellite NSCs to contribute to the final central mass budget. This structure lets us test whether hierarchical GC and NSC assembly can deliver enough black-hole mass to overmassive high-redshift systems without imposing seeds directly by a halo-mass threshold.

The goal of this work is to determine whether GC and NSC assembly can provide a viable origin for high-redshift overmassive black holes, including the black-hole-hosting subset of LRDs. The model is not intended to predict LRD colours, line profiles, obscuration, or radiative-transfer signatures. Instead, it addresses the underlying mass-assembly problem: whether dense clusters can form IMBH seeds early, whether those seeds can be delivered efficiently to galactic centres, and whether the resulting central black holes can become massive compared with their young host galaxies. In this sense, LRDs with confirmed or likely AGN components provide an observational motivation and possible target population, while the model predicts the compact-stellar-system contribution to their black-hole masses. This paper is organised as follows. Sect.~\ref{sect:intro} introduces the observational motivation and the compact-stellar-system seed channel. Sect.~\ref{sect:method} gives a concise overview of the GC formation, sampling, orbital evolution, IMBH-seeding, final-status, and ex-situ NSC merger prescriptions, while the detailed equations and sampling rules are collected in the appendices. The appendices also provide the cosmological relations, galaxy scaling relations, and auxiliary statistical and photometric relations used by the semi-analytic calculation.

\section{Method} \label{sect:method}

The NSC model follows the ex-situ channel in which GCs form during hierarchical galaxy assembly, orbit in an evolving host potential, and contribute stellar mass to the nuclear region through stellar evolution, tidal stripping, direct tidal tearing, and dynamical-friction (DF)-driven sinking. The model is based on the semi-analytic GC framework of \cite{Choksi+2018GC} and \cite{Gao+2024GC}, extended with the \cite{Chen&Gnedin2023GC} baryon-retention limit, the \cite{Chen&Gnedin2024aGC} mass--metallicity relation, the \cite{Fragione+2019GC} local tidal-stripping prescription, and a formation-time IMBH seed module based on the FROST-CLUSTERS runaway-collision fits \citep{Rantala+2026FROST}. This section describes the physical model used for the NSC runs analysed in this work.

\begin{figure*}[htbp!] \centering
\includegraphics[width=\linewidth]{fig/schematic_diagram.png}
\caption{Overview of the semi-analytic GC--NSC--BH pathway. GC formation is painted onto cosmological halo merger trees, with halo growth, gas supply, and baryon retention setting the cluster mass budget. Individual GCs are sampled and assigned birth radii, metallicities, and IMBH seeds; their stellar and orbital evolution is then followed through stellar evolution, tidal stripping, direct tidal tearing, and dynamical friction. Deposited stars build the compact NSC component, while IMBHs that reach the central $1~{\rm pc}$ sink are added to the nuclear BH state, which grows under the central Eddington prescription; satellite branches are evolved recursively and incorporated into their descendants. \czkong{Write corresponding full caption here in the paper style (mimicing other paper figures).}}
\label{fig:schematic_diagram}
\end{figure*}

\subsection{GC formation, sampling, and evolution} \label{subsect:GCformation}

We paint GC formation onto fixed halo merger trees and follow every retained branch. Rapid specific halo growth triggers a formation event; the host stellar-mass, gas, and baryon-retention prescriptions then determine the total cluster mass formed in that event. Individual clusters are sampled from a Schechter initial cluster mass function, assigned formation metallicities, placed in the appropriate main-progenitor or satellite environment, and given birth radii from the corresponding spatial prescription. An IMBH seed is assigned once, at cluster birth, from the cluster mass, half-mass surface density, and metallicity using the runaway-collision fit of \cite{Rantala+2026FROST}.

The clusters are subsequently evolved in a time-dependent NFW plus S{\'e}rsic background that also includes previously deposited GC material. Stellar evolution, continuous tidal stripping, direct tidal tearing, and dynamical friction jointly determine the surviving bound cluster, deposited stellar component, and any wandering IMBH. Objects that reach the fixed $1~{\rm pc}$ sink deposit their remaining stellar component in the innermost radial bin, while an embedded or wandering IMBH enters the central black-hole channel. Appendix~\ref{app:gc_formation} gives the formation trigger and event mass budget; Appendix~\ref{app:gc_sampling} gives the mass, spatial, radii, metallicity, and IMBH sampling rules; and Appendix~\ref{app:evolution_equations} gives the orbital-evolution, mass-loss, disruption, time-stepping, and deposition equations.

\subsection{Eddington-limited nuclear BH growth} \label{subsect:Eddington}

The fiducial run applies Eddington growth only to the stored central BH state. Between two central-entry, branch-import, or output events, the current central mass is advanced in cosmic time according to
\begin{align}
M_\bullet(t + \Delta t) & = M_\bullet(t) \exp\left[\dfrac{f_{\rm Edd}\Delta t}{t_{\rm Edd}}\right], \\
t_{\rm Edd} & = \dfrac{\epsilon_{\rm r}\sigma_{\rm T}c}{4\pi G m_{\rm p}(1 - \epsilon_{\rm r})},
\label{eq:eddington_growth}
\end{align}
where $f_{\rm Edd}$ is the dimensionless input Eddington ratio and $\epsilon_{\rm r}=0.1$ is the radiative efficiency. We use $f_{\rm Edd}=0.3$ for the fiducial output. The central state receives the mass of an IMBH when its host GC or its naked wandering remnant crosses the fixed $1~{\rm pc}$ sink; an intact GC also deposits its remaining stellar mass into the innermost stellar bin. IMBHs that remain in GCs and non-central wandering IMBHs do not accrete. This prescription has no physical upper cap: the $10^{12}~M_\odot$ warning threshold recorded by the run is a diagnostic of excessive exponential growth, not a mass ceiling or a physical prediction.

\subsection{Final Fate of GCs} \label{subsect:stat}

At the end of each evolution segment, the solver writes one final state for every input GC\@. Table~\ref{tab:final_status} summarises the mutually exclusive status flags used in the final-GC catalogue.

\begin{deluxetable*}{c p{0.20\textwidth} p{0.62\textwidth}}
\tablecaption{Final GC and IMBH status flags.\label{tab:final_status}}
\tablehead{\colhead{Status} & \colhead{Fate} & \colhead{Bookkeeping}}
$1$ & Surviving bound GC & A bound stellar cluster remains outside the fixed $1~{\rm pc}$ sink and contributes its surviving stellar mass to the final-GC catalogue. Stellar mass stripped earlier remains in the deposited-mass profile. \\
$-1$ & Stellar exhaustion & The bound stellar component is fully removed outside the sink. The object leaves no surviving GC and no central black-hole entry, but its lost stellar mass remains in the deposited-mass profile. \\
$-2$ & Direct tidal tearing & A non-IMBH-hosting cluster is tidally destroyed outside the sink. Its remaining stellar mass is deposited at the disruption radius. \\
$-3$ & Sunk bound GC & A bound GC reaches the fixed $1~{\rm pc}$ sink before disruption. Its remaining stellar mass is deposited into the innermost radial bin, and any embedded IMBH enters the central black-hole channel. \\
$-4$ & Non-central wandering IMBH & An IMBH-hosting cluster loses its stellar envelope before reaching the sink, leaving a naked IMBH outside the nuclear region at the final redshift. The stripped stellar envelope is deposited where it is removed. \\
$-5$ & Sunk wandering IMBH & A naked wandering IMBH reaches the fixed $1~{\rm pc}$ sink. The IMBH enters the central black-hole channel, with no stellar mass deposited at the entry event.
\tablecomments{The published stellar NSC mass, $M_{\rm NSC}$, is sampled from the cumulative evolved stellar-mass deposition profile inside $R_{\rm NSC}=6~{\rm pc}$. The fixed $1~{\rm pc}$ sink is therefore the innermost deposit bin and the central-BH entry boundary, not the full public NSC aperture.}
\end{deluxetable*}

For IMBH-hosting clusters, loss of the stellar envelope converts the object into a naked wandering IMBH\@. The deposited-mass profile records stellar material removed by stellar evolution, continuous tidal stripping, direct tidal tearing, and sink entry. The stellar NSC mass reported in the halo summaries is then the evolved deposited stellar mass enclosed by $R_{\rm NSC}=6~{\rm pc}$, so stellar material need not be carried intact to $r_{\rm sink}$ in order to contribute to $M_{\rm NSC}$. By contrast, the central black-hole mass includes only IMBHs that cross the fixed $1~{\rm pc}$ sink, either while still embedded in a bound GC or as naked wandering remnants.

\subsection{Halo Mergers \& NSCs} \label{subsect:merge}

The merger treatment follows satellite branches explicitly before their material is handed to the descendant branch. A GC born on a satellite branch is assigned a local S{\'e}rsic radius inside that satellite, using the spatial-sampling procedure in Appendix~\ref{app:gc_sampling} and equation~(\ref{eq:sersic_density}). It is then evolved with the same stellar-evolution, tidal-stripping, tidal-tearing, and dynamical-friction solver described in Appendix~\ref{app:evolution_equations} and equations~(\ref{eq:stellar_evolution})--(\ref{eq:dynamical_friction}), but using the satellite branch's own halo history. The branch is followed until the fixed-tree descendant connection releases it into its recipient branch, or until the terminal branch redshift when no direct descendant row is present.

If a bound GC or a wandering IMBH crosses the fixed $1~{\rm pc}$ sink before the branch merger, its IMBH contribution is stored in the satellite's central black-hole state. The corresponding stellar contribution is retained through the deposited-mass profile: the public $M_{\rm NSC}$ is the evolved deposited stellar mass enclosed within $R_{\rm NSC}=6~{\rm pc}$, including the fixed $0$--$1~{\rm pc}$ bin. At the branch merger, the satellite central black-hole state and its fixed inner-bin deposited stellar component are imported into the recipient branch as compact nuclear material. Surviving non-central GCs and non-central wandering IMBHs are released into the recipient halo at $0.5R_{\rm vir}$ and continue evolving there \citep{Gao+2024GC}. The procedure is recursive, so compact nuclear material assembled in lower-level satellites can pass through intermediate branches before reaching the main progenitor.

\subsection{Cosmological Simulations} \label{subsect:data}

The fiducial catalogue uses $N_{\rm S}=2$ and $f_{\rm Edd}=0.3$ \czkong{No need to mention ex-situ mode 2.} for the calculations presented here \czkong{Not mention run name.}. Its cosmological input contains 2,409 $z=0$ target haloes: 2,105 TNG50-1-Dark haloes selected from the full TNG50 box of side $51.67~{\rm cMpc}$ \czkong{This is actually the full TNG50 box, similar to the TNG100 case.} with $10^{11}<M_{\rm h}\leq10^{13}~M_\odot$, and 304 TNG100-1-Dark haloes selected from the full box with $M_{\rm h}>10^{13}~M_\odot$. The data-retrieval and tree-validation workflow \czkong{Not mention \texttt{dl} dir.} downloads the full SubLink trees for these targets, converts them into the corrected fixed-tree format, and validates the branch ordering and mass-growth invariants before the semi-analytic calculation.

The post-processing combines the formation catalogue, final GC/IMBH states, radial deposition history, redshift-resolved halo summaries, and simulation-provenance tables. \czkong{Not mention any exact file name.} The public $M_{\rm NSC}$ is the evolved deposited stellar mass inside $6~{\rm pc}$, including the innermost $0$--$1~{\rm pc}$ bin. For the QSO1 diagnostics, the deposition profile is matched to the stored $z=7$ sample time, interpolated on a common radial grid, and combined with the central BH mass to obtain the model velocity curve; the UV calculation uses only the GC-origin stellar contributors that had reached the centre before QSO1 and an FSPS--MIST/Chabrier $1500~\AA$ calibration. For abundance plots, TNG50 objects have unit weight in the full TNG50-box reference volume $V_{\rm ref}=(51.67~{\rm cMpc})^3=1.3795\times10^5~{\rm cMpc}^3$, \czkong{This is also a TNG50 full box.} while every TNG100 full-box object is assigned the weight $V_{\rm ref}/V_{\rm TNG100}=0.10164$ before the weighted counts are divided by $V_{\rm ref}$.

\section{Results} \label{sect:result}

The results below use the TNG50--TNG100 catalogue described above. We focus on the central-BH--host relation (Sect.~\ref{subsect:Mbh-Mstar}), the redshift-resolved BH mass functions (Sect.\ref{subsect:BHMF}), and the QSO1 rotation-curve and UV diagnostics (Sect.\ref{subsect:case}). These comparisons test whether GC formation and inspiral can produce a compact, stellar-dominated nucleus that contains a delivered central BH at high redshift.

\begin{comment}
\subsection{Properties of GCs at $z = 0$ \& Comparison with Observations} \label{subsect:check}

\begin{figure}[htbp!] \centering
\includegraphics[width=\linewidth]{fig/Mgc~Mhalo.pdf}
\caption{Mgc vs Mhalo.}
\label{fig:Mgc~Mhalo}
\end{figure}

As a first validation of the underlying GC population, we compare the final GC-system mass with the present-day halo mass. Over the high-mass interval sampled by the fiducial run, the model median follows the approximately linear $M_{\rm GC}$--$M_{\rm h}$ relation inferred for local GC systems and remains close to both the Choksi+2018 model relation and the Harris et al.\ observational calibration. The interquartile range is narrow compared with the full dynamic range of the relation, but it is not negligible; this scatter is expected because the selected haloes have different merger histories, GC formation episodes, metallicities, and subsequent disruption histories.

\begin{figure}[htbp!] \centering
\includegraphics[width=\linewidth]{fig/Age~Metal.pdf}
\caption{Age vs Metal.}
\label{fig:Age~Metal}
\end{figure}

The age--metallicity comparison provides a complementary check on the formation history. The model contours occupy the same broad region as the Galactic and LMC GC measurements: the most metal-poor clusters are predominantly old, while the metal-richer population extends to younger ages and shows larger dispersion. This behaviour is qualitatively consistent with the Choksi-style formation picture used in \citet{Gao+2024GC}, in which early low-metallicity halo growth builds the old metal-poor population and later enrichment shifts subsequent cluster formation to higher metallicity. The agreement should be interpreted as a population-level validation of the GC formation and metallicity prescriptions, rather than a one-to-one fit to the observed Milky Way or LMC samples.

\subsection{The $M_{\rm NSC}$-$M_\star$ Relation} \label{subsect:Mnsc-Mstar}

\begin{figure*}[htbp!] \centering
\includegraphics[width=\linewidth]{fig/Fig.12_Mnsc-Mstar.pdf}
\caption{Mnsc vs Mstar.}
\label{fig:Mnsc~Mstar}
\end{figure*}

The deposited stellar component also falls in the observed NSC scaling regime. In the $M_{\rm NSC}$-$M_\star$ plane, the model points overlap the Neumayer+2020 compilation and the fitted relation for the dynamically measured subset, especially at the massive-galaxy end covered by the present halo sample. The model fit uses the 6 pc deposited-mass proxy, so it should be compared with observed NSC masses as the contribution from inspiralling and disrupted clusters, not as a complete prediction for all nuclear stellar mass.

This distinction is important for interpreting deviations from the observational cloud. The calculation does not include in-situ nuclear star formation, gas inflow, or star formation triggered after cluster deposition. Therefore, agreement with the observed scaling supports the idea that GC inspiral can provide a substantial ex-situ NSC component, while any deficit relative to individual high-mass or late-type systems is not by itself a failure of the ex-situ channel. The declining $M_{\rm NSC} / M_\star$ trend in the right-hand panel is likewise consistent with the observed tendency for NSCs to become a smaller fraction of their host stellar mass in massive galaxies \citep{Neumayer+2020NSC}.

\subsection{The $M_\bullet$-$M_{\rm NSC}$ Relation} \label{subsect:Mbh-Mnsc}

\begin{figure}[htbp!] \centering
\includegraphics[width=\linewidth]{fig/Mbh-Mnsc.pdf}
\caption{$M_{\rm NSC}$-$M_\star$.}
\label{fig:Mbh-Mnsc}
\end{figure}

The model naturally produces systems in which a central BH and an NSC coexist. This is expected observationally: local galaxies can host both components, and the ratio $M_\bullet/M_{\rm NSC}$ spans orders of magnitude \citep{Neumayer+2020}. In the model, the relation is not imposed as an empirical scaling law. The NSC mass is the evolved stellar mass deposited inside $6~{\rm pc}$, whereas the BH mass counts only IMBHs that reach the fixed $1~{\rm pc}$ sink and the subsequent accretion of the central BH state.

The plotted sequence therefore reflects delivery efficiency and accretion history. At high redshift, many haloes have already assembled $M_{\rm NSC}\sim10^7$--$10^9~M_\odot$ but retain comparatively small central BHs, because only a subset of IMBH-bearing clusters have reached the sink. Towards lower redshift, additional inspiral events and central-only accretion move the same population upward in $M_\bullet$ at nearly fixed or slowly growing $M_{\rm NSC}$. By $z=0$, the fiducial catalogue spans $M_\bullet\simeq3.4\times10^7$--$1.6\times10^9~M_\odot$ and $M_{\rm NSC}\simeq4.7\times10^7$--$1.6\times10^9~M_\odot$, illustrating that the model can reach BH-to-NSC mass ratios of order unity in massive descendants.
\end{comment}

\subsection{The $M_\bullet$-$M_\star$ Relation} \label{subsect:Mbh-Mstar}

\begin{figure}[htbp!] \centering
\includegraphics[width=\linewidth]{fig/Fig.01_Mbh-Mstar.pdf}
\caption{Central black-hole mass as a function of the inferred host stellar mass. Coloured solid curves and shaded regions show the binned mean and standard deviation of all finite central-BH entries at each output redshift; dashed curves and their lighter bands retain only entries with $M_\bullet>100~M_\odot$. The colour bar gives redshift. Symbols show the compiled high-redshift observations, the green relation and band show the local Reines--Volonteri relation and its scatter, and dashed green lines mark fixed $M_\bullet/M_\star$ ratios. The host stellar mass is inferred from the main-progenitor halo mass through the adopted SMHM relation. \czkong{Write corresponding full caption here in the paper style (mimicing other paper figures).}}
\label{fig:Mbh-Mstar}
\end{figure}

We first compare the central BH mass with the stellar mass inferred for the same redshift. The stellar mass in this diagnostic is not measured from simulated star particles; it is obtained by applying the project stellar mass-halo mass (SMHM) helper to the main-progenitor halo mass at each output redshift. The comparison is therefore a consistency check between the cluster-seeded central BH channel and an empirical host-galaxy mass scale.

The fiducial tracks occupy the low-mass extrapolation of the local \citet{ReinesVolonteri2015} relation at early times and then move upward as central BH accretion proceeds. At fixed redshift, the model produces a broad range of $M_\bullet/M_\star$, with some high-redshift systems lying above the local relation for their inferred stellar mass. Overall, the model generally matches the observations across the plotted mass and redshift range and reproduces the broad increase of $M_\bullet/M_\star$ towards lower inferred host mass. The remaining scatter reflects the stochastic delivery of IMBHs, the central-only growth prescription, and the fact that the stellar mass on the horizontal axis is an empirical SMHM estimate rather than a directly measured stellar-particle mass.

\subsection{Redshift-resolved BH Mass Function / Number Density} \label{subsect:BHMF}

\begin{figure}[htbp!] \centering
\includegraphics[width=\linewidth]{fig/Fig.03_BHMF.pdf}
\caption{Redshift-resolved black-hole mass function. Coloured solid curves show the number density of positive central BH masses in the halo summaries at the output redshifts, with colour indicating redshift. The grey dash-dotted curve shows positive IMBH masses in the final inventory that remain outside the nucleus. Red and blue curves with shaded regions reproduce the linear and NSC reference mass functions used for comparison. All model densities use the TNG50 reference volume, with TNG100 full-box objects weighted by the ratio of the TNG50 to TNG100 volumes. The Nuclear curves are evaluated at their labelled redshifts, whereas the Satellite curve is a final-inventory diagnostic.}
\label{fig:BHMF}
\end{figure}

For the redshift-resolved BH mass function, we use the component split implemented by the plotting pipeline. The Nuclear component contains positive central-BH masses in the halo-summary rows at each output redshift. The Satellite component contains positive IMBH masses in the final inventory that remain outside the fixed $1~{\rm pc}$ sink, either in a bound GC that survives outside the nucleus or as a naked wandering IMBH whose stellar envelope was lost outside the nucleus. A bound GC or a wandering IMBH that reaches the sink contributes to the central channel instead; clusters whose stellar component is exhausted or tidally torn without a surviving IMBH do not contribute a BH entry. No separate Ejected component is included in this catalogue.

Using the volume convention defined in Sect.~\ref{subsect:data}, the resulting $n_{\rm BH}$ values are weighted catalogue abundances \czkong{This has been mentioned in Sect2.5. Be consistent with previous writting and be brief here since it's repeated.}. In the current output, the Nuclear inventory has 2,409 positive entries at $z=0$, 1,644 at $z=7$, and 534 at $z=10$. The final Satellite inventory contains 3,786,785 positive IMBH entries, with masses from $1.00\times10^2$ to $3.05\times10^3~M_\odot$ before volume weighting. The Nuclear curves are evaluated separately at each output redshift, whereas the Satellite curve is a final-inventory measurement, so the two components should not be interpreted as a single-epoch census. In Figure~\ref{fig:BHMF}, the Nuclear curves remain broadly at the $10^{-3}~{\rm cMpc}^{-3}$ level across the intermediate-mass range at several output redshifts and develop a declining high-mass tail towards later times, while the Satellite curve is concentrated at low masses because it is a final-inventory count. In addition, the uncapped central Eddington prescription can generate late-time central masses above the physically useful plotting range; such warning values are retained for bookkeeping but are not treated as physical SMBH predictions. Figure~\ref{fig:BHMF} is therefore an abundance diagnostic of this mixed simulation sample and its post-processing convention, rather than a complete cosmological BH mass function. \czkong{Need some addtional words to describe the data in the figure.}

\begin{figure}[htbp!] \centering
\includegraphics[width=\linewidth]{fig/Fig.06_BHMF2.pdf}
\caption{Differential black-hole mass function comparison. Coloured curves show the redshift-resolved Nuclear component, the grey dash-dotted curve shows the final Satellite inventory, and the black curve is their $z=0$ sum. The model is plotted as $\log_{10}\Phi$ per dex and uses the TNG50 reference volume with the stated TNG100 volume weights. Symbols and error bars show the compiled high-redshift observational points labelled in the legend; coloured lines and shaded envelopes show the heavy- and light-seed Eddington-limited and super-Eddington reference models. The colour bar gives the output redshift for the Nuclear curves.}
\label{fig:BHMF2}
\end{figure}

Figure~\ref{fig:BHMF2} presents the same model population in the differential form $\Phi(M_{\rm BH})$ per dex, allowing a direct comparison with the compiled high-redshift measurements and the heavy- and light-seed reference tracks. The coloured model curves are the positive central-BH populations sampled at each output redshift, the grey dash-dotted curve is the final population of non-central IMBHs, and the black curve is the exact sum of the $z=0$ Nuclear and Satellite components. The model curves use the same TNG50 reference volume and TNG100 weighting as Figure~\ref{fig:BHMF}; the different presentation changes the binning and logarithmic ordinate, not the underlying catalogue. Across much of the observed $M_{\rm BH}\sim10^6$--$10^8~M_\odot$ interval, the redshift-resolved Nuclear curves overlap the range spanned by the compiled measurements and their uncertainties, while retaining the observed decline in abundance towards larger black-hole masses. The comparison is therefore a broad population-level consistency check, rather than a fit to any individual mass estimate. \czkong{Need some addtional words to describe the data in the figure. Emphasize how the model mataches the obs well.}

\subsection{Case Study: Is Abell 2744-QSO1 a ``naked'' SMBH?} \label{subsect:case}

\begin{figure}[htbp!] \centering
\includegraphics[width=\linewidth]{fig/Fig.02_RotationCurve.pdf}
\caption{Rotation-curve comparison for A2744-QSO1. Symbols show the observed resolved-kinematics, spectroastrometric, and fine-spectroastrometric velocity measurements. The black curve is the QSO1 point-mass comparison with $\log_{10}(M_{\rm BH}/M_\odot)=6.75$, and the grey dash-dotted curve is the Milky-Way-like NSC comparison. The green band encloses the 16th--84th percentiles of the eligible model velocity profiles at $z\simeq7$, while the solid and dashed green curves show their median and mean. The red curve is the equal-weight Keplerian-plus-UV selected halo 221. Negative radii mirror the model profiles to display both sides of the central position.}
\label{fig:RotationCurve}
\end{figure}

The QSO1 comparison is designed to test whether the compact signal attributed to a ``naked'' BH could instead be produced by a compact NSC containing a much smaller central BH. \czkong{Not because this, actually because we want to challenge the idea that QSO1 is a naked BH. We want to prove that it may also be a compact NSC.} We compare the observed rotation-curve points and point-mass curve with the deposited-stellar-plus-central-BH velocity profiles of the model haloes at $z=7$. With equal weights for the Keplerian and UV terms, the lowest finite joint score is $1.244$ and selects halo \texttt{221}; its Keplerian and UV terms are $1.214$ and $0.271$, respectively. The selected line-of-sight velocity profile reaches $48.6~{\rm km~s^{-1}}$ at $12.5~{\rm pc}$, $24.3~{\rm km~s^{-1}}$ at $50~{\rm pc}$, and $14.0~{\rm km~s^{-1}}$ at $150~{\rm pc}$.

At $z=7$, halo 221 contains an evolved deposited stellar mass of $6.72\times10^6~M_\odot$ inside $6~{\rm pc}$ and a central BH mass of $2.12\times10^4~M_\odot$. The enclosed deposited stellar mass is approximately $6.85\times10^6~M_\odot$ across the $12.5$--$150~{\rm pc}$ range, so the central BH contributes only about $0.3$ per cent of the enclosed mass. The BH is included in the velocity calculation, but the predicted amplitude is therefore dominated by the compact stellar component. This result is consistent with a nuclear star cluster assembled by GC inspiral, with IMBH delivery supplying a central BH within the cluster. It supports an embedded-BH interpretation of the compact QSO1 signal, although this particular fiducial halo is not a BH-dominated dynamical match to the observed source.

\begin{figure}[htbp!] \centering
\includegraphics[width=\linewidth]{fig/Fig.05_UVmag.pdf}
\caption{Pure-stellar UV-aperture diagnostic for the QSO1-selected halo 221 at $z=7$. The blue points show the rest-frame $1500~\AA$ absolute AB magnitude of the GC-origin stellar component as the aperture increases from $1$ to $7~{\rm pc}$; the dotted horizontal line marks the QSO1 reference value $M_{\rm UV}=-15.60$. The calculation uses the FSPS--MIST/Chabrier calibration and excludes AGN, nebular, dust, and in-situ young-stellar emission.}
\label{fig:UVmag}
\end{figure}

The UV-aperture diagnostic provides an independent check on the compact stellar interpretation. For halo 221, the GC-origin stellar component has $M_{\rm UV}\simeq-15.46$ at every aperture from $1$ to $7~{\rm pc}$ because essentially all of the selected deposited profile lies inside the innermost bin. This is only $0.14$ mag fainter than the QSO1 reference value, $M_{\rm UV}=-15.60$. The calculation includes the stellar light from the selected GC-origin population only, so it demonstrates that a compact stellar component can account for a large fraction of the reference UV luminosity; it does not determine the contribution from an active central BH, nebular emission, dust, or young in-situ stars. Together with the rotation curve, this supports an interpretation in which a central BH is embedded in a compact NSC assembled by GC inspiral and IMBH delivery.

\begin{figure}[htbp!] \centering
\includegraphics[width=\linewidth]{fig/Fig.04_BHmasses.pdf}
\caption{Black-hole mass comparison for A2744-QSO1. The rows show the compiled QSO1 estimates grouped as virial, scattering, bolometric, and direct estimates; the shaded vertical band marks the MOKA3D estimate and its $1\sigma$ uncertainty. The green plus marker is the model evolved deposited stellar mass inside $6~{\rm pc}$ for the selected halo 221, and the orange cross is its central BH mass at $z=7$. The annotation gives the model pure-stellar UV magnitude and the QSO1 reference magnitude. The model markers describe the compact GC-delivered component and are not a fit to the much larger QSO1 dynamical BH estimate.}
\label{fig:masses}
\end{figure}

For the same joint Keplerian-plus-UV-selected halo, the model gives $M_{\rm NSC}(<6~{\rm pc})=6.72\times10^6~M_\odot$ ($\log_{10} M_{\rm NSC}/M_\odot=6.83$) and $M_\bullet=2.12\times10^4~M_\odot$ ($\log_{10} M_\bullet/M_\odot=4.33$) at $z=7$. The deposited stellar component is therefore about $3.2\times10^2$ times more massive than the central BH in this model case. The $6~{\rm pc}$ pure-stellar UV estimate used in the joint selection is $M_{\rm UV}=-15.46$, close to the QSO1 reference value. These masses summarise a compact GC-delivered nuclear component containing a central BH, rather than a naked BH; the model point should not be read as reproducing the much larger BH masses inferred from the QSO1 dynamical estimates. The comparison instead demonstrates that the observed compact signal can be compatible with an NSC assembled by GC inspiral and a central BH grown from delivered IMBHs.

\begin{figure}[htbp!] \centering
\includegraphics[width=\linewidth]{fig/Fig.09_distr.pdf}
\caption{Halo-mass distributions for the tracked main-progenitor branches. Each coloured step histogram shows the number of available haloes per $0.5$ dex bin in $M_{\rm h}(z)$ at one output redshift, with the colour bar giving redshift. Coloured dashed vertical lines trace the halo-mass history of the selected QSO1-matching halo 221. The dashed track is one selected halo trajectory and is not the distribution of all ten lowest-scoring candidates.}
\label{fig:masses}
\end{figure}

The QSO1 score selects a narrow central-mass window rather than the most massive haloes. Among the 1,586 $z=7$ candidates with finite equal-weight Keplerian-plus-UV scores, the ten lowest-scoring objects are all TNG50 haloes with $9.470\leq\log_{10}(M_{\rm h}/M_\odot)\leq10.209$ and a median $\log_{10}(M_{\rm h}/M_\odot)=9.806$. Their median deposited NSC mass is $7.18\times10^6~M_\odot$, their median initially formed stellar mass inside $6~{\rm pc}$ is $9.44\times10^6~M_\odot$, and their median GC-origin $M_{\rm UV}$ is $-15.47$, close to the QSO1 reference; the median central BH mass is only $1.73\times10^4~M_\odot$. The result is not explained by halo mass alone: among 835 other TNG50 haloes in the same mass interval, the median joint score is $6.297$ and the median $6~{\rm pc}$ initially formed stellar mass is $3.61\times10^6~M_\odot$, compared with $1.386$ and $9.44\times10^6~M_\odot$ for the ten best haloes. The median pre-$z=7$ branch-merger count is 2 for the ten best haloes and 1 for the same-mass control sample (mean values 2.4 and 2.55), so the matching is controlled by the stochastic GC formation, orbital evolution, survival, and central deposition, together with the UV age--metallicity population, rather than by merger count alone. The higher-mass comparison group has a median merger count of 63 but a median Keplerian term of $43.1$, illustrating that excessive central mass can be as unsuccessful as insufficient mass.

To estimate the abundance of such environments, we apply the descriptive dual-match window to all finite-score $z=7$ candidates: $S_{\rm Kep}\leq1.5$ and $|S_{\rm UV}|\leq1$, equivalently $-16.10\leq M_{\rm UV}\leq-15.10$. This selects 21 haloes, all from TNG50, while none of the 304 TNG100 candidates enters the window. With the TNG50 reference volume and the TNG100 weight defined above,
\begin{align}
n_{\rm good}^{\rm selected\ tree}(z = 7) & = \dfrac{21 + 0.10164 \times 0}{1.3795 \times 10^5~{\rm cMpc}^{3}} \\
& = 1.52 \times 10^{-4}~{\rm cMpc}^{-3},
\end{align}
with a Poisson-only uncertainty of $0.33\times10^{-4}~{\rm cMpc}^{-3}$. Varying the descriptive window over the tested thresholds gives values between $0.58\times10^{-4}$ and $2.90\times10^{-4}~{\rm cMpc}^{-3}$; this range represents criterion sensitivity rather than a statistical confidence interval. The fiducial value is a selected-tree density of QSO1-like halo environments that match both the compact stellar kinematics and the GC-origin stellar UV; it is not a number density of observed QSOs or of haloes already hosting a QSO-mass SMBH. The 823 candidates without finite joint scores, cosmic variance, model systematics, and the completeness of the fixed descendant-tree selection are not included in this estimate. \czkong{Based on latest research results in \texttt{report\_paper.md}, \texttt{report\_paper2.md}, and \texttt{report\_paper3.md}, update this para and add the number density estimation results here.}

\section{Discussion \& Caveats} \label{sect:discuss}

%\subsection{Variation of Resolution} \label{subsect:resolution}

\subsection{Variation of Timestep} \label{subsect:timestep}

\subsection{Variation of Scaling} \label{subsect:f_h}

\subsection{Caveats, Limitations, \& Future Works} \label{subsect:limit}

\section{Conclusions \& Summary} \label{sect:summary}

%% Please use the acknowledgment and contribution environments. This will 
%% be anonomyized when the "anonymous" style option is used. 
\begin{acknowledgments}
We thank all the people that have made this AASTeX what it is today.  This
includes but not limited to Bob Hanisch, Chris Biemesderfer, Lee Brotzman,
Pierre Landau, Arthur Ogawa, Maxim Markevitch, Alexey Vikhlinin and Amy
Hendrickson. Also special thanks to David Hogg and Daniel Foreman-Mackey
for the new {\tt\string modern} style design. Considerable help was provided via bug
reports and hacks from numerous people including Patricio Cubillos, Alex
Drlica-Wagner, Sean Lake, Michele Bannister, Peter Williams, Jonathan
Gagne, Arthur Adams, Nicholas Wogan, Aaron Pearlman, Jeff Mangum, Mark Durre, Joel Ong, and Stephen Thorp.
\end{acknowledgments}

\begin{contribution}
%%This section gives authors the space to recognize author contributions. The text inside this environment is NOT counted towards the total word quanta. At a minimum, manuscripts are expected to include this text:

All authors contributed equally to the Terra Mater collaboration.

%% But authors are expected to provide more specific details, e.g. 
%%
%%SC was responsible for writing and submitting the manuscript.
%%WWM came up with the initial research concept and edited the manuscript.
%%OTS obtained the funding and edited the manuscript.
%%EBF provided the formal analysis and validation. He also edited the manuscript.
%%GEH Supervised the undergraduates, wrote the software and administers the project github and Zenodo repositories.
%%
%% Authors can use the Contributor Role Taxonomy (CRediT) at
%% https://credit.niso.org
%% for ideas on how write a good statement tailored to their needs.

\end{contribution}

%% To help institutions obtain information on the effectiveness of their 
%% telescopes the AAS Journals has created a group of keywords for telescope 
%% facilities.
%
%% Following the acknowledgments section, use the following syntax and the
%% \facility{} or \facilities{} macros to list the keywords of facilities used 
%% in the research for the paper.  Each keyword is check against the master 
%% list during copy editing.  Individual instruments can be provided in 
%% parentheses, after the keyword, but they are not verified.
\facilities{HST(STIS), Swift(XRT and UVOT), AAVSO, CTIO:1.3m, CTIO:1.5m, CXO}

%% Similar to \facility{}, there is the optional \software command to allow 
%% authors a place to specify which programs were used during the creation of 
%% the manuscript. Authors should list each code and include either a
%% citation or url to the code inside ()s when available.
\software{astropy \citep{2013A&A...558A..33A,2018AJ....156..123A,2022ApJ...935..167A},  
          Cloudy \citep{2013RMxAA..49..137F}, 
          Source Extractor \citep{1996A&AS..117..393B}
          }

%% Appendix material should be preceded with a single \appendix command.
%% There should be a \section command for each appendix. Mark appendix
%% subsections with the same markup you use in the main body of the paper.
%%
%% Each Appendix (indicated with \section) will be lettered A, B, C, etc.
%% The equation counter will reset when it encounters the \appendix
%% command and will number appendix equations (A1), (A2), etc. The
%% Figure and Table counter will not reset.

\appendix

\section{Cosmology and Halo Quantities}
\label{app:cosmology_equations}

This appendix reviews the cosmological equations used by the semi-analytic model to convert between redshift, cosmic time, halo virial quantities, and the analytic background potential. Throughout these calculations we adopt a standard spatially flat $\Lambda$CDM cosmology and neglect radiation, which is subdominant for the halo-evolution redshifts considered here. The expansion function is
\begin{align}
\Omega_{\Lambda ,0} & = 1 - \Omega_{\rm m, 0} ,\\
E(z) & = \dfrac{H (z)}{H_0} = \sqrt{\Omega_{\rm m, 0} (1 + z)^3 + \Omega_{\Lambda, 0}} .
\end{align}
Here $E(z)$ is the dimensionless Hubble parameter; $H_0$ is the Hubble constant; $H(z)$ is the Hubble parameter at redshift $z$; $\Omega_{\rm m,0}$ is the present-day matter density parameter; and $\Omega_{\Lambda,0}$ is the present-day dark-energy density parameter. The current implementation uses a single cosmology for both the formation catalogues and the analytic orbital-evolution background. We adopt $\Omega_{\rm m,0}=0.307$, $\Omega_{\Lambda,0}=0.693$, and $H_0=67.97\,{\rm km\,s^{-1}\,Mpc^{-1}}$, matching the \cite{DESIcollaboration2025} flat-$\Lambda$CDM constraints used in the code. 

The evolving matter density parameter is
\begin{align}
\Omega_{\rm m}(z) &=
\dfrac{\Omega_{\rm m,0}(1+z)^3}
{\Omega_{\Lambda,0}+\Omega_{\rm m,0}(1+z)^3}.
\end{align}

The cosmic age $t$ at redshift $z$ is
\begin{align}
t(z) = \dfrac{2}{3 H_0 \sqrt{\Omega_{\Lambda,0}}} \operatorname{asinh} \left[ \sqrt{\dfrac{\Omega_{\Lambda,0}}{\Omega_{\rm m,0}}} (1 + z)^{-3/2} \right] .
\end{align}
The inverse age-to-redshift conversion is
\begin{align}
z(t) &= \left[ \dfrac{\sqrt{\Omega_{\Lambda, 0} / \Omega_{\rm m, 0}}}{\sinh \left( 3 H_0 \sqrt{\Omega_{\Lambda,0}} t / 2 \right)} \right]^{2/3} - 1 .
\end{align}

The virial overdensity relative to the critical density follows the \cite{Bryan&Norman1998DM} form
\begin{align}
\Delta_{\rm c} (z) & = 18 \pi^2 + 82 x - 39 x^2 , \\
x &= \Omega_{\rm m}(z)-1 .
\end{align}
The critical density and the corresponding virial radius are
\begin{align}
\rho_{\rm crit} (z) & = \dfrac{3 H^2 (z)}{8 \pi G} , \\
R_{\rm v} & = \sqrt[3]{ \dfrac{3 M_{\rm h}}{4\pi \Delta_{\rm c} (z) \rho_{\rm crit} (z)} } ,
\end{align}
where $G$ is the gravitational constant and $M_{\rm h}$ is halo mass.

The virial velocity is
\begin{align}
v_{\rm v} & = \sqrt{\dfrac{G M_{\rm h}}{R_{\rm v}}} .
\end{align}

For the dark-matter part of the analytic background, the halo concentration and scale radius are
\begin{align}
c &= 9.354
\left(\dfrac{M_{\rm v}h}{10^{12}~M_\odot}\right)^{-0.094},\\
R_{\rm s} &= \dfrac{R_{\rm vir}}{c}.
\end{align}
The radial derivative of the NFW potential is
\begin{align}
\dfrac{\mathrm{d} \Phi_{\rm NFW}}{\mathrm{d} r} = M_{\rm v} \left[ \dfrac{\ln (1 + r / R_{\rm s})}{r^2} - \dfrac{1}{r (r + R_{\rm s})} \right] ,
\end{align}
where $M_{\rm v}$ is the virial mass in the mass unit used by the potential calculation, $R_{\rm s}$ is the NFW scale radius, and $r$ is galactocentric radius.

\section{Galaxy Scaling Relations} \label{appendix:galaxy}

The stellar mass-halo mass(SMHM) relation is written as
\begin{align}
\log_{10} M_{\star} & = \log_{10} (\epsilon M_1) + f \left( \log_{10} \dfrac{M_{\rm h}}{M_1} \right) - f(0) + \xi , \\
f (x) & = - \log_{10} (10^{\alpha x} + 1) + \dfrac{\delta [\log_{10} (1 + e^x)]^\gamma}{1 + e^{10^{-x}}} .
\end{align}
Here $M_{\star}$ is stellar mass, $M_{\rm h}$ is halo mass, $M_1$ is the characteristic halo mass, $\epsilon$ is the stellar-to-halo normalisation, $\alpha$, $\delta$, and $\gamma$ control the low-mass slope, high-mass transition, and curvature, and $\xi$ is the lognormal scatter. The redshift-dependent quantities are
\begin{align}
a & = \dfrac{1}{1 + z} , \\
\nu & = e^{- 4 a^2},\\
\log_{10} \epsilon & = - 1.777 - 0.006 (a - 1) \nu -0.119 (a - 1) , \\
\log_{10} M_1 & = 11.514 + \left[ - 1.793 (a - 1) - 0.251 z \right] \nu , \\
\alpha & = -1.412 + 0.731 (a - 1) \nu , \\
\delta
&= 3.508+\left[2.608(a-1)-0.043z\right]\nu,\\
\gamma
&= 0.316+\left[1.319(a-1)+0.279z\right]\nu,\\
\xi &\sim {\cal N}\left[0,\,0.218-0.023(a-1)\right].
\end{align}

Along a merger-tree branch, stellar mass is grown self-consistently using
\begin{align}
M_{\star,i} &=
M_{\star,{\rm prog}}
+\left[
M_{\star,{\rm SMHM}}(M_{{\rm h},i},z_i)
-M_{\star,{\rm SMHM}}(M_{{\rm h},{\rm prog}},z_{\rm prog})
\right]10^{\xi_i}.
\end{align}
Here the subscript ``prog'' denotes the first progenitor, and $M_{\star,{\rm SMHM}}$ is the median stellar mass from the SMHM relation.

The cold-gas mass is expressed through $\eta=M_{\rm g}/M_{\star}$:
\begin{align}
\log_{10}\eta &=
-0.45 - n_M(\log_{10}M_{\star}-9) + Z_{\eta}(z) + \xi_{\rm g},\\
n_M &=
\begin{cases}
0.19, & M_{\star}<10^9~M_\odot,\\
0.33, & M_{\star}\geq10^9~M_\odot,
\end{cases}\\
\xi_{\rm g} &\sim {\cal N}(0,0.3).
\end{align}
The redshift term is
\begin{align}
Z_{\eta}(z) &=
\begin{cases}
(3-t_{\rm dep})\log_{10}\left(\dfrac{1+z}{3}\right)
+(3-t_{\rm dep})\log_{10}3, & z<2,\\
(1.7-t_{\rm dep})\log_{10}\left(\dfrac{1+z}{3}\right)
+(3-t_{\rm dep})\log_{10}3, & 2\leq z<3,\\
(1.7-t_{\rm dep})\log_{10}\left(\dfrac{4}{3}\right)
+(3-t_{\rm dep})\log_{10}3, & z\geq3,
\end{cases}\\
t_{\rm dep} &= 0.3.
\end{align}
The resulting gas mass is $M_{\rm g}=\eta M_{\star}$ before applying the baryon-retention ceiling.

The retained baryon fraction is limited by
\begin{align}
f_{\star} &= \dfrac{M_{\star}}{f_{\rm b}M_{\rm h}},\\
f_{\rm g} &= \dfrac{M_{\rm g}}{f_{\rm b}M_{\rm h}},\\
f_{\rm g} &\rightarrow
\begin{cases}
0, & f_{\star}\geq f_{\rm in},\\
f_{\rm in}-f_{\star}, & f_{\star}<f_{\rm in}\ {\rm and}\ f_{\star}+f_{\rm g}>f_{\rm in},\\
f_{\rm g}, & f_{\star}+f_{\rm g}\leq f_{\rm in},
\end{cases}
\end{align}
where $f_{\rm b}$ is the cosmic baryon fraction, $f_{\star}$ and $f_{\rm g}$ are stellar and gas fractions normalised to $f_{\rm b}M_{\rm h}$, and $f_{\rm in}$ is the accreted baryon fraction.

The smooth filtering function used for baryon retention is
\begin{align}
s(x,y) &= \left[1+\left(2^{y/3}-1\right)x^y\right]^{-3/y}.
\end{align}
The Chen--Gnedin retention model is
\begin{align}
f_{\rm in} &= s\left(\dfrac{M_{\rm ch}}{M_{\rm h}},2\right),\\
M_{\rm ch}(z) &=
\dfrac{1.69\times10^{10}e^{-0.63z}}
{1+e^{(z/\beta)^{\Gamma}}}~M_\odot,\\
\beta &= z_{\rm rei}
\left[\ln\left(1.82\times10^3 e^{-0.63z_{\rm rei}}\right)-1\right]^{-1/\Gamma},\\
z_{\rm rei} &= 6,\qquad \Gamma=15.
\end{align}
Here $M_{\rm ch}$ is the characteristic halo mass at which reionisation suppresses baryon accretion.

The production runs use the \cite{Chen&Gnedin2024aGC} mass-metallicity relation,
\begin{align}
[\mathrm{Fe} / \mathrm{H}] (M_\star, z) & = 0.3 \log_{10} \left( \dfrac{M_\star}{10^9~M_\odot} \right) - \log_{10} (1 + z) - 0.5 .
\label{eq:mass-metal}
\end{align}
This relation is capped at $[\mathrm{Fe}/\mathrm{H}]=0.3$. The cluster metallicity is
\begin{align}
[\mathrm{Fe}/\mathrm{H}]_{\rm GC}
&=[\mathrm{Fe}/\mathrm{H}]_{\rm gal}+\xi_m,\\
\xi_m &\sim {\cal N}(0,0.3),
\end{align}
and the solar-scaled metal fraction entering the IMBH fit is
\begin{align}
\dfrac{Z}{Z_{\odot}} &= 10^{[\mathrm{Fe}/\mathrm{H}]_{\rm GC}}.
\label{eq:metallicity_transfer}
\end{align}

\section{GC Formation}
\label{app:gc_formation}
\label{app:cluster_initial_equations}

GC formation is painted onto fixed halo merger trees and is followed along all retained branches. A formation event is triggered when the specific halo growth rate exceeds $p_3$,
\begin{align}
\dfrac{M_{{\rm h},i}/M_{{\rm h},{\rm prog}}-1}
{t_i-t_{\rm prog}} & > p_3 ,
\label{eq:formation_trigger}
\end{align}
where $t_i$ and $t_{\rm prog}$ are the cosmic ages of the current halo and its progenitor. The parameter $p_3$ sets the formation-trigger threshold, while $p_2$ sets the cluster-formation efficiency. The host stellar-mass, gas, and baryon-retention relations are given in Appendix~\ref{appendix:galaxy}. The total cluster mass assigned to one formation event is
\begin{align}
M_{\rm GC, tot} & = 3 \times 10^{-5} p_2 \dfrac{M_{\rm g}}{f_{\rm b}},
\label{eq:gc_total_mass}
\end{align}
Events with $M_{\rm GC, tot}<10^5~M_\odot$ do not form a cluster in the current implementation.

\section{GC Sampling}
\label{app:gc_sampling}
\label{subsect:gc_sampling}

After the event mass budget is fixed, the model samples the individual clusters in a sequence of mass, spatial, radius, metallicity, and IMBH assignments. The sampling choices are described below; the subsequent orbital evolution is given in Appendix~\ref{app:evolution_equations}.

\subsection{Mass sampling}
\label{subsubsect:mass_sampling}

Individual cluster masses are drawn from a Schechter initial cluster mass function,
\begin{align}
\dfrac{\mathrm{d}N}{\mathrm{d}M} & = A_{\rm c}M^{-2}e^{-M/M_{\rm c}},
\label{eq:cimf}
\end{align}
where $A_{\rm c}$ is the normalisation and $M_{\rm c}$ is the exponential cut-off mass. We use $M_{\min}=10^5~M_\odot$; the fiducial runs use $M_{\rm c}=10^7~M_\odot$. The upper incomplete gamma function is
\begin{align}
\Gamma(s,x) & = \int_x^{\infty} u^{s-1}e^{-u}\,\mathrm{d}u .
\end{align}
The deterministic maximum-mass relation is obtained from the condition that one cluster lies above $M_{\max}$,
\begin{align}
1 & = A_{\rm c}M_{\rm c}^{-1}\Gamma\left(-1,\dfrac{M_{\max}}{M_{\rm c}}\right),\\
M_{\rm GC,tot} & =
\dfrac{M_{\rm c}
\left[
\Gamma\left(0,\dfrac{M_{\min}}{M_{\rm c}}\right)
-\Gamma\left(0,\dfrac{M_{\max}}{M_{\rm c}}\right)
\right]}
{\Gamma\left(-1,\dfrac{M_{\max}}{M_{\rm c}}\right)}.
\end{align}
For sampling between $M_{\min}$ and $M_{\max}$, the cumulative probability is
\begin{align}
r(<M) & =
\dfrac{
\Gamma\left(-1,\dfrac{M_{\min}}{M_{\rm c}}\right)
-\Gamma\left(-1,\dfrac{M}{M_{\rm c}}\right)}
{
\Gamma\left(-1,\dfrac{M_{\min}}{M_{\rm c}}\right)
-\Gamma\left(-1,\dfrac{M_{\max}}{M_{\rm c}}\right)}.
\end{align}
Sampling continues until the event mass budget is exhausted. If the final draw exceeds the remaining budget, its mass is set to the remaining mass. The cluster list is shuffled before radii are assigned, so the deterministically reserved maximum-mass cluster is not always placed at the smallest S{\'e}rsic radius.

\subsection{Spatial sampling}
\label{subsubsect:spatial_sampling}

Main-progenitor-branch clusters are treated as in-situ objects. Clusters formed on satellite branches are first assigned positions in their own local host, before any later halo merger. If a satellite cluster survives outside the nuclear region until the branch merges, it is released into the recipient halo at $0.5R_{\rm vir}$ and continues evolving in the new host potential. Thus, the branch on which a cluster forms determines whether its first radius is local to the main progenitor or to a satellite system.

\subsection{Radii sampling}
\label{subsubsect:radii_sampling}

For main-progenitor and locally sampled satellite clusters, the birth radius is drawn from the deprojected S{\'e}rsic profile used by \cite{Gao+2024GC},
\begin{align}
\rho(r) & = \rho_0 \left(\dfrac{r}{R_{\rm e}}\right)^{-p}
e^{-b\left(r / R_{\rm e}\right)^{1 / N_{\rm S}}},
\label{eq:sersic_density}
\end{align}
where
\begin{align}
p & = 1 - \dfrac{0.6097}{N_{\rm S}} + \dfrac{0.05563}{N_{\rm S}^2}, \\
b & = 2 N_{\rm S} - \dfrac{1}{3} + \dfrac{0.009876}{N_{\rm S}}.
\end{align}
The effective radius follows the halo-spin scale,
\begin{align}
R_{\rm e} & = \dfrac{j_{\rm sp}h}{20H(z)R_{\rm vir}},
\end{align}
where $j_{\rm sp}$ is the specific angular-momentum norm of the halo, $h=H_0/(100~{\rm km~s^{-1}~Mpc^{-1}})$, $H(z)$ is the Hubble parameter, and $R_{\rm vir}$ is the virial radius. The enclosed-mass fraction is inverted through the regularised incomplete gamma function,
\begin{align}
r_{\rm GC} & = R_{\rm e}
\left[\dfrac{P^{-1}\{N_{\rm S}(3-p),u\}}{b}\right]^{N_{\rm S}},
\end{align}
where $u$ is the assigned enclosed mass fraction within the formation event. Accreted clusters that are released after a branch merger use the representative outer radius
\begin{align}
R_{\rm acc} & = 0.5R_{\rm vir}.
\end{align}
In the analytic background, the corresponding spin parameter and effective radius are
\begin{align}
\lambda & = \dfrac{J}{\left(2GR_{\rm vir}M_{\rm h}\right)^{1/2}},\\
R_{\rm e} & = \dfrac{\lambda R_{\rm vir}}{\sqrt{2}},
\end{align}
where $J$ is the halo angular-momentum norm in the same unit system as $G$, $R_{\rm vir}$, and $M_{\rm h}$.

\subsection{Metallicity assignment}
\label{subsubsect:metallicity_assignment}

The host-galaxy metallicity and the $0.3$ dex cluster-to-cluster scatter are given in Appendix~\ref{appendix:galaxy}. The resulting $[\mathrm{Fe}/\mathrm{H}]_{\rm GC}$ is used for stellar-evolution bookkeeping and is converted to the solar-scaled abundance entering the IMBH fit through equation~(\ref{eq:metallicity_transfer}).

\subsection{IMBH sampling}
\label{append:IMBH}
\label{subsect:IMBH}
\label{subsubsect:runaway_imbh}

The IMBH module assigns a seed once, at GC birth, using the runaway-collision prescription of \cite{Rantala+2026FROST}. Dense, metal-poor clusters can build very massive stars through stellar collisions before those products collapse into IMBHs. The cluster structural scale is the three-dimensional half-mass radius
\begin{align}
r_{\rm h} & =
\dfrac{f_{\rm h}R_4}{1.3}
\left(\dfrac{M_{\rm cl}}{10^4~M_\odot}\right)^\beta,
\label{eq:imbh_radius}
\end{align}
with $f_{\rm h}=0.125$, $R_4=2.365~{\rm pc}$, and $\beta=0.180$. Assuming a Plummer profile, the projected half-mass radius and half-mass surface density are
\begin{align}
R_{\rm h} & = \dfrac{r_{\rm h}}{1.305},\\
\Sigma_{\rm h} & = \dfrac{M_{\rm cl}}{2\pi R_{\rm h}^2}.
\label{eq:imbh_sigma}
\end{align}

The fit is calibrated over $10^{-2}\leq Z/Z_{\odot}\leq1$. The implementation evaluates the supplied positive metallicity ratio in its corresponding piecewise interval and issues a warning when it lies outside this range; it does not clip the input.
%The metallicity entering the IMBH fit is restricted to the calibrated interval:
%\begin{align}
%\left( \dfrac{Z}{Z_{\odot}} \right)_{\rm IMBH} & =\min \left[ 1, ~\max \left( 10^{-2}, ~\dfrac{Z}{Z_{\odot}} \right) \right] .
%\end{align}
For each metallicity interval, every fit coefficient $Y$ is a linear function of $\log_{10}(Z/Z_{\odot})$,
\begin{align}
Y & = k_Y\log_{10}\left(\dfrac{Z}{Z_{\odot}}\right)+b_Y,
\end{align}
where $Y$ denotes one of the tabulated coefficients and $(k_Y,b_Y)$ is selected from the corresponding metallicity segment. Let
\begin{align}
Q & = \log_{10}\left(\dfrac{\Sigma_{\rm h}}
{M_\odot~{\rm pc}^{-2}}\right).
\end{align}
In the calibrated-density regime, the IMBH mass fit is
\begin{align}
M_{\bullet,9} & =
\begin{cases}
AQ+BQ^2+C, & \Sigma_{\rm h}\geq\Sigma_{\rm crit},\\
0, & \Sigma_{\rm h}<\Sigma_{\rm crit},
\end{cases}
\end{align}
where the critical density is
\begin{align}
\log_{10}\left(\dfrac{\Sigma_{\rm crit}}
{M_\odot~{\rm pc}^{-2}}\right) & = S.
\end{align}
For the high-density branch,
\begin{align}
M_{\bullet,10} & = DQ+E.
\end{align}
The selected seed mass is
\begin{align}
M_{\bullet} & =
\begin{cases}
M_{\bullet,10}, & Q\geq5.22,\\
M_{\bullet,9}, & Q<5.22,
\end{cases}\\
M_{\bullet} & \rightarrow
\begin{cases}
M_{\bullet}, & M_{\bullet}\geq100~M_\odot,\\
0, & M_{\bullet}<100~M_\odot.
\end{cases}
\end{align}
Here $A$, $B$, $C$, $D$, $E$, and $S$ are the metallicity-dependent fitted coefficients. The assigned IMBH mass and the birth structural diagnostics $(r_{\rm h},\Sigma_{\rm h})$ are stored with the GC and are not updated by later stellar evolution or tidal stripping.
\czkong{Add the coef table here.}

\section{Cluster Evolution}
\label{app:evolution_equations}
\label{subsect:gc_evolution}

The evolution calculation follows every formed GC in an analytic, evolving background from its formation time to the requested final redshift. The halo mass and spin are interpolated along the relevant merger-tree branch. The dark-matter component is modelled with an NFW profile, the stellar component with a S{\'e}rsic background, and previously deposited model GCs contribute to the enclosed mass used at later timesteps. The background is updated in coarse time blocks, while the fastest-changing active object is advanced with adaptive mass-loss and orbital-decay steps.

\subsection{Evolving background and orbital decay}

The stellar background mass enclosed by radius $r$ is
\begin{align}
M_{\star}(<r) &=
M_{\star}
P\left[n(3-p),\,b\left(\dfrac{r}{R_{\rm e}}\right)^{1/n}\right].
\end{align}
The circular velocity of the smooth background is
\begin{align}
v_{\rm c,bg}^2(r) &=
r\dfrac{\mathrm{d}\Phi_{\rm DM}}{\mathrm{d}r}
+\dfrac{M_{\star}(<r)}{r},
\end{align}
and the associated mean background density is
\begin{align}
\rho_{\rm bg}(r) &=
\dfrac{v_{\rm c,bg}^2(r)}{(4\pi/3)r^2}.
\end{align}
The deposited cluster material already enclosed within $r$ contributes
\begin{align}
\rho_{\rm GC}(<r) &=
\dfrac{M_{\rm GC}(<r)}{(4\pi/3)r^3},
\end{align}
with unit conversions absorbed into the numerical constants of the evolution equations. The circular velocity used for orbital decay is
\begin{align}
v_{\rm c} (r) = \sqrt{\dfrac{4 \pi }{3} G \left( \rho_{\rm bkgd} + \rho_{\rm GC} \right) r^2}
\end{align}

The dynamical-friction inspiral rate is
\begin{align}
\dfrac{\mathrm{d}r}{\mathrm{d}t} &=
-\dfrac{M_{\rm GC,5}}{0.45\,r\,v_{\rm c}},
\label{eq:dynamical_friction}
\end{align}
where $M_{\rm GC,5}=M_{\rm GC}/(10^5~M_\odot)$, $r$ is in kpc, and $v_{\rm c}$ is in km s$^{-1}$. The fourth-order Runge--Kutta update evaluates the positive inward step function $D(r)=M_{\rm GC,5}/(0.45rv_{\rm c})$ as
\begin{align}
k_1 &= \Delta t\,D(r),\\
k_2 &= \Delta t\,D(r-k_1/2),\\
k_3 &= \Delta t\,D(r-k_2/2),\\
k_4 &= \Delta t\,D(r-k_3),\\
\Delta r &= \dfrac{k_1+2k_2+2k_3+k_4}{6}.
\end{align}
The radius is then updated as $r\rightarrow r-\Delta r$.

\subsection{Mass loss and disruption}

For the local-orbit tidal-stripping model \citep{Fragione+2019GC},
\begin{align}
P &= 100 \dfrac{r}{v_{\rm c}} , \\
\dfrac{\mathrm{d} M_{\rm GC}}{\mathrm{d} t} & = \dfrac{2^{2/3} M_{\rm GC}^{1/3}}{10 P} .
\label{eq:tidal_stripping}
\end{align}
Here $P$ is a normalised orbital period. The alternative fixed-period disruption law is \citep{Choksi+2018GC}
\begin{align}
t_{\rm tid} & = 5~{\rm Gyr} \left( \dfrac{M_{\rm GC}}{2 \times 10^5~M_\odot} \right)^{2/3} , \\
t_{\rm iso} & = 17~{\rm Gyr} \left( \dfrac{M_{\rm GC}}{2 \times 10^5~M_\odot} \right) , \\
\dfrac{\mathrm{d} M_{\rm GC}}{\mathrm{d} t} & = - \dfrac{M_{\rm GC}}{\min(t_{\rm tid}, t_{\rm iso})} .
\end{align}

The analytic formation-stage survival estimate applies
\begin{align}
m_0 &= \dfrac{M_0}{2\times10^5~M_\odot},\\
t_{\rm tid,0} &= 10^{10}~{\rm yr}\,m_0^{2/3}P_0,\\
\nu_{\rm tid,0} &= \dfrac{1}{t_{\rm tid,0}},\\
k(t) &= 1-\dfrac{2}{3}\nu_{\rm tid,0}(t-t_{\rm form}),\\
m_{\rm tid}(t) &= m_0 k^{3/2}.
\end{align}
If weak-field evaporation is included, the transition mass is
\begin{align}
m_{\rm iso} &= \left(\dfrac{P_0}{1.7}\right)^3,
\end{align}
and, when $m_{\rm tid}<m_{\rm iso}$,
\begin{align}
t_{\rm iso,start} &=
t_{\rm form}
+\dfrac{3\left[1-(m_{\rm iso}/m_0)^{2/3}\right]}
{2\nu_{\rm tid,0}},\\
m_{\rm f} &=
m_{\rm iso}
-\dfrac{t-t_{\rm iso,start}}{1.7\times10^{10}~{\rm yr}}.
\end{align}

The stellar-evolution mass-loss fraction is represented as
\begin{align}
f_{\rm sw}(t) &=
\max\left[0,\,-\dfrac{x^2}{100}+0.288x-1.42\right],\\
x &= \log_{10}\left(\dfrac{t}{\rm yr}\right),
\label{eq:stellar_evolution}
\end{align}
For tabulated metallicity-dependent stellar evolution, the solar and subsolar tables are used directly at $[\mathrm{Fe}/\mathrm{H}]\geq0$ and $[\mathrm{Fe}/\mathrm{H}]\leq -0.69$. Between them, the stellar-evolution time grid is interpolated as
\begin{align}
t_{\rm int} &=
t_{\rm sub}
+\dfrac{t_{\odot}-t_{\rm sub}}{0.69}
\left([\mathrm{Fe}/\mathrm{H}]+0.69\right).
\end{align}
The final mass after analytic survival filtering is
\begin{align}
M_{\rm f} &=
2\times10^5~M_\odot\,m_{\rm dyn}
\left[1-f_{\rm lost}(t-t_{\rm form},[\mathrm{Fe}/\mathrm{H}])\right],
\end{align}
where $m_{\rm dyn}$ is either $m_{\rm tid}$ or the weak-field value $m_{\rm f}$, and $f_{\rm lost}$ is the metallicity-dependent stellar-evolution loss fraction.

\subsection{Stellar evolution and direct tidal tearing}

During the time-dependent evolution, the stellar-wind loss over one step is
\begin{align}
\Delta f_{\rm sw} &=
f_{\rm sw}(t+\Delta t-t_{\rm form})
-f_{\rm sw}(t-t_{\rm form}),\\
\Delta M_{\rm sw,bound} &=
M_{\rm init}\dfrac{M_{\rm bound}}
{M_{\rm bound}+M_{\rm dep}}\Delta f_{\rm sw},\\
\Delta M_{\rm sw,dep} &=
M_{\rm init}\dfrac{M_{\rm dep}}
{M_{\rm bound}+M_{\rm dep}}\Delta f_{\rm sw},\\
\Delta M_{\rm tid} &=
\left|\dfrac{\mathrm{d}M_{\rm GC}}{\mathrm{d}t}\right|\Delta t.
\end{align}
Here $M_{\rm bound}$ is the current bound cluster mass, $M_{\rm dep}$ is the mass already deposited by the same cluster, and $M_{\rm init}$ is its initial mass.

The cluster half-mass density used for direct tidal tearing is approximated as
\begin{align}
\rho_{\rm h} & =
\begin{cases}
10^3, & M < 10^5~M_\odot,\\
10^3 \left(\dfrac{M}{10^5~M_\odot}\right)^2, & 10^5~M_\odot \le M \le 10^6~M_\odot,\\
10^5, & M > 10^6~M_\odot,
\end{cases}
\label{eq:cluster_density}
\end{align}
in the internal density units of the evolution solver. A cluster is torn apart when its half-mass density falls below the sum of the analytic background and deposited-GC densities,

\begin{align}
\rho_{\rm h} < \rho_{\rm bkgd}+\rho_{\rm GC} .
\end{align}

\subsection{Wandering IMBH remnants}
\label{subsubsect:wandered_imbh}

An IMBH-hosting GC can lose its stellar envelope before reaching the centre. This happens either when continuous stellar and tidal mass loss reduces the bound stellar mass down to the IMBH mass, or when direct tidal tearing removes the remaining stellar component. The model then deposits the residual stellar mass into the local radial bin and converts the object into a wandering IMBH\@. A wandering IMBH no longer undergoes stellar evolution or tidal stripping, but it continues to experience dynamical friction using its own mass. If it reaches $r_{\rm sink}$, it is counted as a sunk wandering IMBH; otherwise it remains as a surviving off-centre remnant at the final redshift.

The adaptive time-step limits are
\begin{align}
\Delta t_M &= t_{\rm s,m}\dfrac{M_{\rm GC}}{|\mathrm{d}M_{\rm GC}/\mathrm{d}t|},\\
\Delta t_r &= t_{\rm s,r}\dfrac{r}{|\mathrm{d}r/\mathrm{d}t|},\\
\Delta t &= \min(\Delta t_M,\Delta t_r,\Delta t_{\max}),
\end{align}
where $t_{\rm s,m}$ and $t_{\rm s,r}$ are user-specified fractional step parameters.

The radial deposition bins are logarithmic:
\begin{align}
b(r) &= 1+\left\lfloor
\dfrac{\log_{10}\max(r,r_{\min})-\log_{10}r_{\min}}
{\log_{10}r_{\max}-\log_{10}r_{\min}}
(N_{\rm bin}-1)
\right\rfloor ,
\end{align}
with the result limited to the interval $1\leq b\leq N_{\rm bin}$.

\section{Auxiliary Statistical and Photometric Relations}
\label{app:auxiliary_equations}

For a power-law probability density $\mathrm{d}N/\mathrm{d}M\propto M^{\alpha}$, inverse-transform sampling gives
\begin{align}
M &= M_{\min}(1-r)^{1/(\alpha+1)}
\end{align}
for an unbounded upper tail, and
\begin{align}
M &=
\left[
r\left(M_{\max}^{\alpha+1}-M_{\min}^{\alpha+1}\right)
+M_{\min}^{\alpha+1}
\right]^{1/(\alpha+1)}
\end{align}
for the bounded interval $M_{\min}\leq M\leq M_{\max}$, where $r$ is a uniform deviate on $[0,1)$.

The Gaussian kernel density estimate used for one-dimensional distributions is
\begin{align}
\widehat{p}(x) &=
\dfrac{1}{\sqrt{2\pi}\,h\,N}
\sum_{i=1}^{N} w_i
e^{-(x-x_i)^2/(2h^2)},
\end{align}
where $h$ is the bandwidth, $x_i$ are the samples, $w_i$ are optional weights, and $N$ is the number of samples.

Linear interpolation on a uniformly spaced grid is
\begin{align}
y(x) &= y_j + q(y_{j+1}-y_j),\\
j &= \left\lfloor \dfrac{x-x_0}{\Delta x}\right\rfloor,\\
q &= \dfrac{x-x_0}{\Delta x}-j.
\end{align}

\section{Rest-Frame UV Magnitude from the Model Stellar Mass}
\label{app:mstar_to_muv}

This appendix describes the conversion between model stellar mass and rest-frame ultraviolet (UV) magnitude used for the QSO1 comparison. The procedure is a forward estimate of the intrinsic stellar light from the compact, GC-origin nuclear component. In the QSO1 UV comparison, the relevant mass is instead the initially formed stellar mass deposited inside a small nuclear aperture,
\begin{align}
M_{\star,{\rm form}}(<R_{\rm UV}),
\end{align}
with $R_{\rm UV}=6~{\rm pc}$ for the current QSO1 selection. This aperture mass is used because the UV calibration table gives the luminosity per unit initially formed stellar mass, not per unit surviving stellar mass.

The calculation is intended to answer a limited question: if the compact nuclear stellar component assembled by GC inspiral were observed at the QSO1 redshift \czkong{$z = 7.04$ in the QSO1 case and $z \approx 7$ in our model.}, what rest-frame $1500~\AA$ absolute magnitude would its old stellar population produce? The default stellar estimate contains only the stellar continuum. It does not include dust attenuation, nebular emission, reprocessed AGN light, intergalactic absorption, or X-ray-binary emission. These omitted components can change the observed UV light, but they are not part of the pure-stellar conversion described here.

For each candidate halo, the model first identifies the GC-origin contributors that have already reached the nuclear region before the QSO1 epoch. If a GC formed at cosmic time $t_{{\rm form},i}$ and the comparison is made at redshift $z_{\rm QSO1}$, the stellar age assigned to that contribution is
\begin{align}
\tau_i &= t(z_{\rm QSO1})-t_{{\rm form},i}.
\end{align}
Equivalently, in terms of lookback time,
\begin{align}
\tau_i &= t_{{\rm lb},{\rm form},i}-t_{\rm lb}(z_{\rm QSO1}).
\end{align}
Each contributor also carries the metallicity assigned at GC formation, written as $[\mathrm{Fe}/\mathrm{H}]_i$. These ages and metallicities are not used to change the total aperture mass directly. Instead, they determine the mass-to-light ratio of the composite old population.

The stellar-population calibration is a tabulated simple-stellar-population relation,
\begin{align}
m_{1500} &= m_{1500}(\tau,[\mathrm{Fe}/\mathrm{H}]),
\end{align}
where $m_{1500}$ is the rest-frame $1500~\AA$ absolute AB magnitude produced by $1~M_\odot$ of initially formed stars. The current table was generated for MIST stellar tracks, a Chabrier initial mass function, and a pure stellar continuum. The tabulated variables are age and metallicity. The interpolation is carried out in $\log_{10}\tau$ and $[\mathrm{Fe}/\mathrm{H}]$, because the UV mass-to-light ratio changes rapidly at young ages and more smoothly in logarithmic age. If a candidate lies outside the calibrated age or metallicity interval, the estimate is evaluated at the nearest tabulated boundary and the number of such boundary evaluations is recorded as a diagnostic. Thus the calculation avoids an uncontrolled extrapolation beyond the stellar-population table.

The individual tabulated magnitudes cannot be averaged directly, because magnitudes are logarithmic. The composite specific luminosity per unit initially formed stellar mass is therefore averaged in luminosity space. If the $i$th contributor has a statistical weight $w_i$ and interpolated magnitude $m_{1500,i}$, define
\begin{align}
q_i &= 10^{-0.4 m_{1500,i}},\\
q_{\rm comp} &=
\dfrac{\sum_i w_i q_i}{\sum_i w_i},\\
m_{1500,{\rm comp}} &=
-2.5\log_{10} q_{\rm comp}.
\end{align}
Here $q_i$ is proportional to the $1500~\AA$ luminosity per unit initially formed stellar mass. The proportionality constant cancels in the weighted average, so this step only determines the effective stellar-population magnitude per solar mass. In practice, the weights trace the relative importance of the sunk GC contributors when setting the age and metallicity mixture. The total luminosity normalisation is still set by the aperture mass $M_{\star,{\rm form}}(<R_{\rm UV})$.

The stellar UV magnitude of the compact nuclear component is then
\begin{align}
M_{\rm UV} & = m_{1500,{\rm comp}} - 2.5\log_{10} \left[ \dfrac{M_{\star,{\rm form}}(<R_{\rm UV})}{M_\odot} \right] .
\label{eq:muv_stellar_from_mstar}
\end{align}
This equation is the central $M_\star$-to-$M_{\rm UV}$ conversion used in the QSO1 comparison. A larger formed stellar mass makes the object brighter by adding more copies of the same composite stellar population. A younger or more metal-poor population can also make the object brighter by lowering $m_{1500,{\rm comp}}$, because such a population emits more UV light per unit formed mass.

For reference, the AB-magnitude conversion used for adding luminosity components is
\begin{align}
L_{\nu}(M_{\rm AB})
&=10^{0.4(51.60-M_{\rm AB})}
~{\rm erg~s^{-1}~Hz^{-1}},\\
M_{\rm AB}(L_{\nu})
&=-2.5\log_{10}
\left(
\dfrac{L_{\nu}}{{\rm erg~s^{-1}~Hz^{-1}}}
\right)
+51.60 .
\end{align}

The QSO1 analogue search uses the resulting magnitude as one part of the candidate score. The observed target value adopted for the comparison is
\begin{align}
M_{{\rm UV},{\rm QSO1}} &= -15.60 ,
\end{align}
with a tolerance of $0.5~{\rm mag}$. The UV residual is therefore
\begin{align}
\Delta_{\rm UV} &=
\dfrac{M_{\rm UV}-M_{{\rm UV},{\rm QSO1}}}{0.5}.
\end{align}
This residual is combined with the rotation-curve residual in the final selection statistic. The UV term should be interpreted as a consistency filter for choosing QSO1-like model nuclei, not as an independent fit to the full observed spectrum.

The main assumptions behind this conversion are therefore explicit. The stellar mass is the initially formed GC-origin mass inside the chosen nuclear aperture; the age and metallicity mixture is inferred from GCs that are already present before the comparison epoch; the stellar population is treated as a pure stellar continuum; luminosities are averaged and added in linear space; and the optional AGN component is added only after the stellar magnitude has been converted to luminosity.

%% For this sample we use BibTeX plus aasjournalv7.bst to generate the
%% the bibliography. The sample7.bib file was populated from ADS. To
%% get the citations to show in the compiled file do the following:
%%
%% pdflatex sample7.tex
%% bibtext sample7
%% pdflatex sample7.tex
%% pdflatex sample7.tex

\bibliography{ref}{}
\bibliographystyle{aasjournalv7}

%% This command is needed to show the entire author+affiliation list when
%% the collaboration and author truncation commands are used.  It has to
%% go at the end of the manuscript.
%\allauthors

%% Include this line if you are using the \added, \replaced, \deleted
%% commands to see a summary list of all changes at the end of the article.
%\listofchanges

\end{CJK*}
\end{document}

% End of file `sample7.tex'.
